\chapter{Conclusão e Próximos Passos}

O projeto consolidado define uma plataforma de pesquisa tecnicamente defensável: radomes geodésicos distribuídos, uma abertura híbrida com antenas Yagi externas para VHF e elementos menores para frequências superiores, recepção vetorial multifaixa, processamento local de eventos, sincronização precisa, calibração ponta a ponta e fusão multiestática. O princípio central não é uma antena universal, mas uma família de antenas mecanicamente integrada e compartilhando uma interface comum de sincronismo, calibração e fusão na face.

Como inovação arquitetural candidata, o projeto combina as Yagis externas com células piramidais receptoras individualmente blindadas e com revestimento absorvedor ainda a selecionar. A hipótese é preservar o caminho de recepção fora do casco e, simultaneamente, reduzir auto-interferência, acoplamento e reradiação pela abertura. Essa possível baixa observabilidade deve ser comprovada por balanço de potência, eficácia de blindagem, padrões OTA e RCS angular; o documento não afirma invisibilidade nem desempenho operacional.

O conjunto de Yagis de faixas cruzadas é uma solução proposta de integração mecânica, ainda não uma inovação demonstrada. Uma Yagi VHF externa maior e uma Yagi UHF menor montada ortogonalmente combinam funções espectrais distintas em uma interface de face, mas permanecem canais independentes de polarização única e não são combinadas para síntese polarimétrica. A polarimetria completa fica reservada a módulos coerentes dual-porta dentro da mesma faixa.

O próximo passo de menor risco é um demonstrador de três nós VHF/UHF com canais de referência e vigilância, seguido por tiles L/S/C e X/Ku/Ka. HF deve avançar como programa próprio de sensores vetoriais e modos característicos da plataforma. A expansão territorial deve aguardar a medição de cobertura, orçamentos de enlace, geometria, falso alarme e estabilidade metrológica.
